\chapter{Formalisation of Lilac}
% Here I present the current progress in the formalisation of Lilac, and how I went about it. I divide it into 2 parts
% Initial work which was mostly measure theoretic, and recent work which is more structured.
% As in any project its best to get to the most uncertain parts of it first, to minimise its risk. I think I have somewhat
% intuitively followed this principle in my formalisation so far, also meaning its full of \texttt{sorry}s.
% However the shape of the formalisation
% of Lilac is reasonably clear and I was able to identify some noteworthy problems which show up in the formalisation.

% The repository can be found on \href{https://github.com/edwin1729/bppl-logic}{GitHub}: It contains most of the syntax and semantics of both the language and program logic for it. 
% These can be regarded as the "data" parts (as opposed to "proof" parts) of the formalisation. There was much consideration (and confusion) in modelling
% both language and logic in a principled manner. We want a representations for free variables and contexts where we don't go out of bounds by construction.
% The approach is inspired by Dependent de Bruijn Indices \cite{cpdt} which also shows how types/propositions and terms in the language/logic are translated consistently
% into lean types and terms repsectively when taking the denotation.

% The goal with defining these, data parts is to be able to, using MoSeL, reimplement the case studies
% provided in the Lilac paper in an ergonomic manner (with tactical support), with the proofs of soundness of the logic itself being left as \texttt{sorry}.
% The "proof" parts are mostly incomplete, with the proof of unitarity of the
% unital element of the KRM being almost complete, being a first milestone (as it involves proof of corollaries of the $\pi$-$\lambda$-theorem)

% As a brief reflection on the difference in formalisation on the maths and computer science sides, in my limited experience,
% I do think in this project (computer science) I do need to think of multiple things at once and papers generally disseminate
% dependencies more clearly without depending as heavily on the reader to morph and modify lemmas as required. However
% the objects in Lilac are so unconventional from the measure theory in mathlib, that it requires inventing some things from
% scratch and quite challenging. This is ideal for my goals of understanding how to use a dependent type theory for formalisation. 
% As such \ref{pi-lambda} talking about formalisation specific issues. But \ref{dependent-store}
% talks about something more conceptual in Lilac! (My pivot towards the aim of a more theoretical contribution is that
% I feel I have gotten enough Knowledge of Lean from formalising Lilac)

% \section{Measure spaces as a dependent type}
% \label{paired-prob-space}

% \section{Parametrisation of the "store" by the "heap" in this separation logic}
% \label{dependent-store}
% The classical separation logic assertions are written like so $s, h \vDash P$, where the $s$ is a variable store and
% $h$ is a heap. In the satisfiability relation, none of $s$, $h$ or $P$ are dependent on another argument. Consider the
% probabilistic separation logic with satisfiability relation of form $\gamma, D, \mathcal{P} \vDash P$, 
% where $\gamma$ is the deterministic variable store, $D$ is the random variable store and $\mathcal{P}$ is the 
% probability space analagous to heap. I claim that $D$ is dependent on $\mathcal{P}$, and that many of the conditions in
% the satisfiability relation are actually true by construction.

% This unspotted dependent typing is a consequence of how measurable functions are usually just denoted $A \overset{m}{\rightarrow} B$
% with the $\sigma$-algebra on A and B that is required being generally obvious from the context (there is usually only one).
% However in our case, we our changing the $sigma$-algebra we're working with all the time when taking independent combinations,
% in the separating conjunct.

% The relevant measurable function here is $D$ the random variable store. Its type is
% $D : [0, 1]^{\mathbb{N}} \overset{m}{\rightarrow} \llbracket \Delta \rrbracket$, where $\llbracket \Delta \rrbracket$ is the 
% denotation of the random variable context, a list of types the random values that can occur in the separation logic proposition
% $P$. As seen in a PSL, $D$ has quite an interesting type as a measurable function which can be composed with an expression
% in a propostion, which is equivalent to substitution for the usual deterministic variables.
% Now no explicit $\sigma$-algebra is provided for the Hilbert cube, the domain of $D$ and the natural candidate for this is
% the $sigma$-algebra for the hilbert cube contained in the probability space $\mathcal{P}$. This interpretation yeilds
% some simplifications in the satisfiability relation of probability specific connectives, exemplified in the \texttt{Own}
% connective. It's semantics is:

% $\gamma, 𝐷, ([0, 1]^{\mathbb{N}}, F, \mu) \vDash own \, 𝐸 \quad iff \quad 𝐸 (\gamma) \circ 𝐷$ is F-measurable
% where $\mathcal{P} = ([0, 1]^{\mathbb{N}}, F, \mu)$ and $F$ is $\mathcal{P}$'s $\sigma$-algebra.
% So $𝐸 (\gamma) \circ 𝐷$ is F-measurable by construction. Ownership as measurability is a key concept of Lilac, and it neatly falls into place.

\section{Syntax and semantics of Appl}
To reason about Appl (see \ref{sec:appl}), we first need to define Appl in Lean. A first attempt at doing this may be inspired by how compilers are written:
define an AST for the syntax of the language, and a typechecker for accepting or rejecting AST terms. Following such an approach would be painful given our goal
of specifying properties of programs. We would always have to guard for the possibility that a program is not well-typed, in the assertion, and would lead to a
proliferation of exception cases. Fortunately we are working in a dependently typed language, which is expressive enough to circumvent the problem by
letting us define a syntax (AST) that is well-typed by construction.

For disambiguation, we refer to types in the \emph{meta language} (in our case Lean) and \emph{object language} (Appl), when it is unclear.

There are several standard ways to define object languages which are well-typed by construction, which all differ in the crucial problem of representing
(object language) variables. This seemingly innocent problem of how to represent variables is one of the biggest issues encountered in this project, and
standard techniques don't quite work for Appl, because of it's non-standard interpretation of (object language) types.

The final (Lean) representation of chosen for Appl follows a technique called Dependent De Brujin Indices
\citep[prefix][see chap 17, Reasoning about programming language syntax]{cpdt}. Discussion on the alternatives are postponed to \ref{}.

First we present the types of Appl along with its denotational semantics, taken from \href{https://github.com/edwin1729/bppl-logic/blob/submission/Bppl/Lilac/Appl.lean}{Bppl.Lilac.Appl}.
\begin{leancode}
inductive Ty where
  | prod : Ty → Ty → Ty
  | bool
  | real
  | exp : ℕ → Ty → Ty -- Tyⁿ
  | index
  | G : Ty → Ty

@[reducible] def Ty.den : Ty → MeasCat
  | prod ty₁ ty₂ => .of (ty₁.den × ty₂.den)
  | bool => .of Bool
  | real => .of ℝ
  | exp n ty => .of (∀ (_ : Fin n), ty.den) -- using MeasurableSpace.pi
  | index => .of ℕ
  -- using `MeasureTheory.ProbabilityMeasure.instMeasurableSpace`
  | G ty => .of (ProbabilityMeasure (ty.den))
\end{leancode}

This mirrors the definition in the paper with no deviations. Note that the denotation of Appl types is \lean|MeasCat|, the category of
MeasurableSpace, rather than \lean|Type| as for most other object languages. It is crucial to annotate \lean|@[reducible]|, which tells the
typechecker to unfold \lean|Ty.den| when checking if two types are definitionally equal. The importance of this will be made clear once we describe
the denotation of terms.

First we describe the 
De Brujin indices 

Here is an excerpt of the definitions of terms, (omitted constructors can be found on Github).

\begin{leancode}
inductive Term : List Ty → Ty → Type
  | var : Member ty ctx → Term ctx ty
  | ret : Term ctx ty → Term ctx (.G ty)
  | bind : Term ctx (.G ty₁) → Term (ty₁ :: ctx) (.G ty₂) → Term ctx (.G ty₂)
  | pair : Term ctx ty₁ → Term ctx ty₂ → Term ctx (.prod ty₁ ty₂)
  | fst : Term ctx (.prod ty₁ ty₂) → Term ctx ty₁
  | snd : Term ctx (.prod ty₁ ty₂) → Term ctx ty₂
  | T : Term ctx .bool
  | F : Term ctx .bool
  | ite : Term ctx .bool → Term ctx ty → Term ctx ty → Term ctx ty
  | flip : (p : ℝ≥0) → (h : p ≤ 1) → Term ctx (.G .bool)
  | r : ℝ → Term ctx .real
  | arith : Arith → Term ctx .real → Term ctx .real → Term ctx .real
  | cmp : Cmp → Term ctx .real → Term ctx .real → Term ctx .bool
  | unif01 : Term ctx (.G .real)
\end{leancode}

\begin{leancode}
/-- Takes the term and variable environment (which is a measurable function)
to give an element of a measurable space -/
@[simp] def Term.den : Term ctx ty → List.TProd (⟪·⟫) ctx -m→ ty.den
  | var mem => ⟨fun env ↦ env.get mem, List.TProd.measurable_get mem⟩
  | ret X => ⟨fun env ↦ probDirac (X.den env), measurable_probDirac.comp X.den.2⟩
      -- MeasureTheory.Measure.measurable_dirac.comp X.den.2⟩
  -- In `bind` we are working exactly with kernels (as shown by the two coercions above)
  -- The `Kernel` api bundles `measurable` property and we need to convert to it and back
  -- to use its deep measure theoretic results for constructing measurable functions
  | @bind _ ty₁ ty₂ M N =>
    letI T := List.TProd (⟪·⟫) ctx
    letI k₁ : Kernel T (⟪ty₁⟫ × T) := (toMK M.den ×ₖ Kernel.id)
    letI k₂ : Kernel (⟪ty₁⟫ × T) ⟪ty₂⟫ := toMK N.den -- need to provide type annotation for this
    toDen (k₂ ∘ₖ k₁)
  | pair M N => ⟨fun env ↦ (M.den env, N.den env), Measurable.prod M.den.2 N.den.2⟩
  | fst M => ⟨fun env ↦ (M.den env).fst, Measurable.fst M.den.2⟩
  | snd M => ⟨fun env ↦ (M.den env).snd, Measurable.snd M.den.2⟩
  | T => ⟨fun env ↦ true, measurable_const⟩
  | F => ⟨fun env ↦ false, measurable_const⟩
  | ite P M N => ⟨fun env ↦ if P.den env then M.den env else N.den env,
      -- the function translating from bool to Prop is `Measurable`
      -- by `measurableSet_singleton true`
      Measurable.ite (P.den.2 (measurableSet_singleton true)) M.den.2 N.den.2⟩
  | flip p hp => ⟨fun _ ↦ ⟨(bernoulli p hp).toMeasure, inferInstance⟩, measurable_const⟩
  | r x => ⟨fun _ ↦ x, measurable_const⟩
  | arith op M N => ⟨fun env ↦ op.den (M.den env, N.den env),
    (op.den.2).comp (M.den.2.prod N.den.2)⟩
  | cmp op M N => ⟨fun env ↦ op.den (M.den env, N.den env),
    (op.den.2).comp (M.den.2.prod N.den.2)⟩
  | unif01 => ⟨fun _ ↦ ⟨unif01_sem, inferInstance⟩, measurable_const⟩
\end{leancode}

interesting because of the
 \lean|theorem foo : True := by trivial|
  Some things that are noteworthy are the fact 

\section{Semantics of Lilac propositions}
\subsection{Reviewing Probabilistic Separation Logics}
More recently there has been a lot of research into probabilistic separation logics: \citep{og_psl}, introduced separation as probabilistic independence.
\citep{lilac} introduced Lilac, which allowed usage of continuous distributions, providing a measure theoretic definition of resources. \citep{basl} introduced reasoning about bayesian
updates (supporting the \texttt{observe} and \texttt{normalize} constructs), thereby providing an axiomatic semantics for the general form of first order
probabilistic programs (ie, BPPL as introduced by \citep{prob-semantics}).

There is a parallel line of work in relational separation logics. One could argue that relational logics are more important in the probabilistic domain, since in bayesian
probabilistic programing (so including \texttt{observe} and \texttt{normalize}), frequently generates probability distributions which have no closed form solution! Sampling
from these can only be done using approximation methods like Variational Inference. The justification for relational logics classical is that relational properties are
expressed more naturally, but could perhaps still be done using unary reasoning, in an inelegant way. In probabilistic programs however, relational reasoning might
in fact be more expressive making it important. Relational reasoning for probabilistic programs (though not using a separation logic) can be found here \cite{relational-logic}.
More recently there has been a novel combination of relational and unary reasoning for probabilistic programs in a separation logic by \citep{bluebell}.

\subsection{Lilac}
\label{lilac}
The following are a few core intuitions behind Lilac. Lilac aims to reason about random variables, which is a higher level of abstraction that we generally
reason in informally (as opposed to descending down into measure theory and lebesgue integrals). Since it's a separation logic, it is also compositional but
along a different axis to compositionality of the denotational semantics. I call the compositionality of probabilistic separation logic a line-by-line compositionality
(as opposed to term-wise compositionality). Each let binding creates a new "line" and we can compose our reasoning about each of line. In the denotational semantics
we would instead reason compositionally about the arguments to the monadic bind operator (an s-finite measure and an s-finite-kernel). The former is a more natural way to reason
about a sequencing operator.

This reasoning style is achieved by defining a resource which gets carried through and composed and split according to the laws of separation logic. This resource or state is
a canonical source of randomness, the Hilbert cube, along with a measure space structure over it. It can be argued that the Hilbert cube could be replaced with something else,
and indeed the definitions are parametric over this choice of measure space for a while. Any distribution can be constructed from the Hilbert cube (which is like a uniform distribution)
using the operations/terms of the language.

\Todo{ownership as measurability}

\Todo{independent combinations of probability spaces}

% It is unclear to me why an axiomatic semantics would need to be provided in terms of measurable maps
% from the Hilbert cube. I wonder whether measure theory is not abstract enough and perhaps synthetic
% measure theory may hold some intuitions/inspiration or simplifications 


\section{Instantiating Iris}
Iris \cite{iris} is a step-indexed modal separation logic partly designed as a framework for embedding other separation logics into,
to take advantage of the Iris Proof Mode (IPM) \cite{ipm}. The Iris proof mode provides state of the art "tactical support", tactic
based interactive theorem proving. More recently the front of the IPM was abstracted from Iris to any separation logic to
create MoSeL \cite{mosel}. We aim to use MoSeL in this project for several reasons
\begin{enumerate}
\item Neither Lilac nor BaSL is step-indexed, so it is much cleaner to instantiate MoSeL rather than Iris
\item Making Lilac compatible with Lilac's interface (instantiating an algebraic structure called CMRA \cite{iris}) is an
unsolved problem (future work section \cite{lilac})
\item MoSeL was completely ported to Lean, shortly before the start of this project! \cite{iris-lean}.
\end{enumerate}

MoSeL presents the interface of a "logic of Bunched Implications" which essentially contains all standard connectives of a first order
logic and the separation logic connectives. The only addition MoSeL adds to this is a persistence modality (which can be instantiated in 
a trivial way in exchange for defeating any automation it might achieve?).

One of the key  generalisations that MoSeL makes over Iris is parametricity over linear and affine logics. Iris is an affine logic
meaning separating conjuncts in a premise can be used 0 or 1 time in proving a goal; they can be thrown away if desired. In a linear
logic however every conjunct must be used exactly once, as in the first separation logic by \citet{og_sl}.
There is naturally quite an emphasis on describing whether a logic is affine or linear, since throwing away unwanted predicates
when it is suitable to do so, is exactly the kind of thing a user wants automation for. In BaSL \cite{basl}, there is a short
discussion stating BaSL is part affine and part linear. Precisely characterising the constructs which are affine and linear or under
what conditions should be explored for better tactical support.

MoSeL like IPM before it provides tactics mirroring the native Lean tactics, such as \texttt{iIntros} and \texttt{IAssumption}

In an ongoing project to port Iris from Rocq to Lean, MoSeL was completely ported to Lean right before the beginning of this project
\cite{iris-lean-impl, iris-lean}. So the formalisation builds on this. The interface provided by MoSeL is quite general. The interface
defines a pre-orderered Set  (PROP, Entails)
\begin{leancode}
class BIBase (PROP : Type u) where
%   ⊢ : PROP → PROP → Prop -- Entails
  pure : Prop → PROP
  ∧ : PROP → PROP → PROP
  → : PROP → PROP → PROP
  * : PROP → PROP → PROP -- separating conjunct
  -∗ : PROP → PROP → PROP -- magic wand
  persistently : PROP → PROP
  ...
\end{leancode}

The type class \texttt{BIBase} defines the interface for the "data" parts of the separation logic, while the axioms that these
need to satisfy (the "proof" parts) are defined in an extension of this type class.  We are defining the separation logic, Lilac or BaSL, in the meta-logic, Lean. 
So \texttt{BIBase} is parametrised by a type \texttt{PROP} which is a proposition in the separation logic (as opposed to \texttt{Prop} in lean).

We draw the reader's attention to the \texttt{Entails} which allows us to provide the semantics of entailment. The soundness of MoSeL's
syntactic entailments are established by the semantics provided when instantiating \texttt{Entails}. This gives us a little flexibility
in how we'd like to instantiate \texttt{PROP}. For simpler logics we could have \texttt{PROP}, directly capture the semantis of that
proposition, giving \texttt{PROP} this shape. 

\begin{leancode}
PROP := Store → Heap → Prop
\end{leancode}

My understanding is that this precisely encodes a relation: the satisfiability relation. An n-ary relation is encoded as a function
from the n parameters to a Proposition. In other words it's a set of n-tuples. (Sets of type A are encoded as $A \rightarrow Prop$).
Notice however that $s, h \vDash P$ should be of the form 

\begin{leancode}
PROP := Store → Heap → PROP → Prop
\end{leancode}

where did the $P$ (of type \texttt{PROP}) go? It's given implicitly through the inductive construction of a term of \texttt{PROP}
through the constructors defined by \texttt{and}, \texttt{sep}, etc. Any term of \texttt{PROP} implicitly contains the tree of
constructors used to create it.

This is standard for the simple example instantiations of MoSeL (such as classical separation logic), but quite convoluted
conceptuall. We may wish to split the representations of the separation logic assertions into syntax and
semantics, which is what is done in non-trivial logics and this projects formalisation. There is a denotation function of type
$Store \rightarrow Heap \rightarrow PROP \rightarrow Prop$ (exactly how the satisfiability relation should look) 
used in the instantiation of Entails; much more natural!

\section{Pi-lambda-theorem as an induction principle}
\label{pi-lambda}
The $\pi$-$\lambda$-theorem is tool in measure theory which can be used for equality of two measures. We use it prove a certain uniqueness of measures theorem,
used to establish that separating conjunction yeilds a PCM (partial commutative monoid) structure.
It is similar to an induction principle on the structure of a $\sigma$-algebra. So we need prove some property C holds in the base and inductive cases
of a $\sigma$-algebra. However there isn't an inductive type underlying this
induction principle! The reason is two fold 1) when generating a measurable set from some base measurable sets, a specific
set can be shown to be measurable by multiple different permutations of application of the axioms. For reference below
is included an inductive defintion generating a $\sigma$-algebra from base measurable seats.

\begin{leancode}
/-- The smallest σ-algebra containing a collection `s` of basic sets -/
inductive GenerateMeasurable (s : Set (Set α)) : Set α → Prop
  | protected basic : ∀ u ∈ s, GenerateMeasurable s u
  | protected empty : GenerateMeasurable s ∅
  | protected compl : ∀ t, GenerateMeasurable s t → GenerateMeasurable s tᶜ
  | protected iUnion : ∀ f : ℕ → Set α, (∀ n, GenerateMeasurable s (f n)) →
      GenerateMeasurable s (⋃ i, f i)
\end{leancode}

2) The $\pi$-$\lambda$-theorem provides a relaxation on our proof obligation. We dont prove that the property C holds for countable unions 
(\texttt{iUnion} above) if it holds for each of the sets we're taking a union over. Instead we only need to do so for countably pairwise disjoint unions.

Being closed under disjoint unions instead of unions, is the difference between a $\lambda$-system compared to a $\sigma$-algebra. This relaxtion is the 
content of the $\pi$-$\lambda$-theorem which we don't go into here.
However there is an alternative "inductive" definition of a $\lambda$ system with the following axioms: 
\begin{enumerate}
\item The system $\mathcal{C}$ containt the base measurable sets we chose
\item be closed under proper difference, for $\forall V, U \in \mathcal{C}, V \subseteq U \rightarrow U \setminus V \in \mathcal{C}$
\item closed under increasing limits, take $i \in \mathbb{N}$ and $U_i \in \mathcal{C}$ then $\bigcup_{i \in \mathbb{N}} U_i \in \mathcal{C}$
\end{enumerate}

Unfortunately this is not the definition of $\lambda$-system in Mathlib, and it doesn't provide the "induction principle" according to this definition.
For the case of probability measures ($\mu$ (universal set) = 1), there is a simple workaround, which is sufficient for Lilac. However BaSL uses unnormalised
measures since it support bayesian conditioning, and so, there may be no simple workaround but reproving this alternative induction principle.

\section{Binder Discipline}
The choice between Dependent De Brujin Indices and Parametric Abstract Higher Order Syntax (PHOAS).
\subsection{Dependent De Brujin Indices}
\begin{itemize}
\item More straightforward to the uninitiated following a little more intuitively from the paper. Important if
we want this is to be an accessible code base for future probabilistic separation logics to follow from.
\item Measurability of maps expressible explicitly
\end{itemize}

\subsection{PHOAS}
\begin{itemize}
\item Notion of substitution inbuilt
\item Example programs written in it are naturally readable. No separate parser and translation from a normal
lambda calculus required. Correctness of the parser doesn't need to be proven... (but further consideration on
this point requires deciding if there will be a separate executable implementation of the APPL syntax)
\end{itemize}

\section{Equality of dependently typed records in Lean}


\section{Choice of Kripke Resource Monoid (KRM) over Unital ordered Resource algebra (ORA) }
ORA's are established as the standard by former Rocq-iris projects, but I couldn't see why not just directly
use the KRM. It was conceptual simplicity traded for some possibly unweildy proofs. The unweildiness turned out to be
restricted to a single file?

\section{Deep vs shallow embedding}

\section{Support for recursion over the data structure even with (shallow?/deep? see SSProve talk) embedding}
(compare with SProp usage in Rocq for ORA? or something...)
Even though the monotonicity property did not require induction, substitution atleast as defined in the paper requires it.

The aim, is to do things as given in the paper, and then think about if substitution could be done away with? If removing
the ability to do recursion over assertions is desirable, etc

\section{Strong abstractions across the proof development}
\subsection{The language}
This is built on the existing abstractions. We use the typeclass \texttt{DenotationalMeas} to decouple the syntax from
the language from the logic. The language only needs to me able to interpreted as measurable maps with a few additional
requirements as given by \texttt{DenotationalMeas}.

\subsection{Resource monoid}
\emph{for the version of Lilac with/without conditioning}

\subsection{Separate treatment of BI connectives and Probability specific connectives}

\section{Metaprogramming and tactical support}

% - The correct type for meaurable spaces because type casting issue, it needed to be a dependent type.

% - Correctly instantiating BI, modelling random variables, "intrinsically type calculus"? 
% - Talk about Lilac, give the intuition that this is the actual line by line semantics